\section{Fibonacci fusion operators}%
\label{sec:asymex_fibonacci}

The Fibonacci fusion category has two simple objects, the vacuum $1$
and $\tau$, with the fusion rule $\tau\otimes\tau=1\oplus\tau$. Its
string-net construction produces matrix product operators whose fusion
rule, fusion tensors, and $F$-symbols are given
in~\cite[Appendix~D.1]{Bultinck2017Anyons}; the review displays the
string-net operator built from the same Fibonacci data
in~\cite[Appendix~A, ``The MPO for the Fibonacci model'']{Cirac2021Matrix},
and the general string-net construction of such operator symmetries
and their intertwiners is in~\cite{Lootens2021MPOStringNet}. The
physical space used here is the two-label space of the golden
chain~\cite{Feiguin2007GoldenChain}, one label $1$ or $\tau$ per site,
rather than the full triple-index alphabet of the categorical
string-net construction. The entries of
the blocks defined below are obtained from the $F$-symbols
of~\cite[Appendix~D.1]{Bultinck2017Anyons}. From these data one obtains
two operators on the chain whose fusion relation and action on periodic
states reduce to finite matrix computations.

The physical space at each site is $\C^2$, with computational labels
$0$ for the vacuum and $1$ for $\tau$. At system size $N\geq1$, write
$\mc H_N=(\C^2)^{\otimes N}$ and use the orthonormal basis
$\ket{x}=\ket{x_0,\ldots,x_{N-1}}$. Site indices are cyclic:
$x_N=x_0$. Throughout, the golden-number convention is
$\varphi=(1+\sqrt5)/2$ and
$\sigma=\varphi^{-1/2}=\sqrt{(\sqrt5-1)/2}$.
In particular, $\sigma^2=\varphi^{-1}$ and
$\sigma^4+\sigma^2=1$.

For an MPO tensor $T^{yx}$, the first physical index is the output
index. Its periodic operator is defined by
\begin{align}
    \bra{y}O_N(T)\ket{x}
     & =\tr(T^{y_0x_0}\cdots T^{y_{N-1}x_{N-1}}).
    \label{eq:asymex_fib_periodic}
\end{align}
For the tensors $B_1$ and $B_\tau$ defined below, write
$O_1=O_N(B_1)$ and $O_\tau=O_N(B_\tau)$ when $N$ is fixed.
Their non-invertible fusion relation is
\begin{align}
    O_\tau O_\tau & =O_1+O_\tau,
    \label{eq:asymex_fib_fusion}
\end{align}
as proved below for every positive system size. The vacuum operator
$O_1$ is a projection onto admissible configurations, not the identity
on $\mc H_N$.

Every entry of the tensors below, and of the changes of bond
coordinates in the proofs, lies in a quartic ring generated by
$\sigma$. The identities of this section are first verified in that
ring and then embedded into $\C$.

\begin{definition}[The ring of golden integers]
    \label{def:asymex_golden_ring}
    \lean{GoldenInt, GoldenInt.sigma, goldenPoly, goldenSigmaReal,
        goldenToComplex, MPSTensor.complexOfGolden, MPSTensor.evalWordGolden,
        MPSTensor.mulGoldenTensor}
    \leanok
    The ring of \emph{golden integers} is $\Z[\sigma]$, where
    $\sigma=\varphi^{-1/2}=\sqrt{(\sqrt5-1)/2}$ is the positive real
    root of $X^4+X^2-1$. It is presented by the four integer
    coordinates of $c_0+c_1\sigma+c_2\sigma^2+c_3\sigma^3$,
    with multiplication reduced using $\sigma^4=1-\sigma^2$.
    It embeds into $\C$ through the quotient
    $\mathbb{Q}[X]/(X^4+X^2-1)$ evaluated at this real root.
    For matrices, the embedding is applied entrywise.
    Word evaluation and the bond-space product of tensors over
    $\Z[\sigma]$ use the same formulas as over $\C$.
\end{definition}

\begin{theorem}[Exact arithmetic in the ring of golden integers]
    \label{thm:asymex_golden_ring}
    % Principal declaration: goldenMulRule.
    \lean{GoldenInt.sigma_quartic, goldenMulRule, goldenSigmaReal_quartic,
        MPSTensor.complexOfGolden_mul, MPSTensor.evalWord_complexOfGolden,
        MPSTensor.mulTensor_complexOfGolden}
    \leanok
    \uses{def:asymex_golden_ring}
    The generator satisfies $\sigma^4+\sigma^2-1=0$ in $\Z[\sigma]$, and so
    does the real number $\sqrt{(\sqrt5-1)/2}$. In any commutative ring
    carrying a root $s$ of $X^4+X^2-1$, the product of two integer
    combinations of $1,s,s^2,s^3$ is the combination prescribed by the
    multiplication of $\Z[\sigma]$. The entrywise embedding of matrices
    over $\Z[\sigma]$ into the complex matrices is multiplicative, and it
    commutes with word evaluation and with the bond-space product of
    tensors.
\end{theorem}

\begin{proof}
    \leanok
    Squaring $\sigma^2=(\sqrt5-1)/2$ gives the quartic for the real root.
    The multiplication rule is the reduction of a product of two
    polynomials of degree at most three modulo $X^4+X^2-1$.
    It is a polynomial identity in their eight integer coefficients.
    The entrywise embedding is induced by a ring homomorphism, so
    expansion of a matrix product proves its multiplicativity.
    Induction on the word proves compatibility with word evaluation.
    Expansion of the bond-space product proves the remaining
    compatibility statement.
\end{proof}

\begin{definition}[The two Fibonacci blocks and the stacked square]
    \label{def:asymex_fib}
    \lean{FibonacciCompression.fibTau, FibonacciCompression.fibOne,
        FibonacciCompression.fibTauMPS, FibonacciCompression.fibOneMPS,
        FibonacciCompression.fibStack, FibonacciCompression.fibTargets}
    \leanok
    \uses{def:asymex_golden_ring, def:mpo_tensor, def:mpdo_operator_product}
    Read a physical letter as the pair $(x',x)$ of outgoing and incoming
    labels, ordered as $(0,0),(0,1),(1,0),(1,1)$. The reduced vertex
    weight is
    \begin{align}
        F & =\begin{pmatrix}
                 0 & \varphi^{-1}   & \varphi^{-1/2} \\
                 1 & 0              & 1              \\
                 1 & \varphi^{-1/2} & -\varphi^{-1}
             \end{pmatrix}.
        \label{eq:asymex_fib_vertex}
    \end{align}
    The $\tau$ tensor $B_\tau$ has bond dimension three and letters
    \begin{align}
        B_\tau^{00} & =0.
        \label{eq:asymex_fib_tau_zero}
    \end{align}
    \begin{align}
        B_\tau^{01} & =\begin{pmatrix}
                           0 & \sigma^2 & \sigma \\0&0&0\\0&0&0
                       \end{pmatrix}.
        \label{eq:asymex_fib_tau_01}
    \end{align}
    \begin{align}
        B_\tau^{10} & =\begin{pmatrix}
                           0 & 0 & 0 \\1&0&1\\0&0&0
                       \end{pmatrix}.
        \label{eq:asymex_fib_tau_10}
    \end{align}
    \begin{align}
        B_\tau^{11} & =\begin{pmatrix}
                           0 & 0 & 0 \\0&0&0\\1&\sigma&-\sigma^2
                       \end{pmatrix}.
        \label{eq:asymex_fib_tau_11}
    \end{align}
    Thus each nonzero letter carries one row of
    (\ref{eq:asymex_fib_vertex}), with the other rows zero.
    The vacuum tensor $B_1$ has bond dimension two and letters
    \begin{align}
        B_1^{00} & =\begin{pmatrix}0&1\\0&0\end{pmatrix}.
        \label{eq:asymex_fib_one_00}
    \end{align}
    \begin{align}
        B_1^{11} & =\begin{pmatrix}0&0\\1&1\end{pmatrix}.
        \label{eq:asymex_fib_one_11}
    \end{align}
    \begin{align}
        B_1^{01} & =B_1^{10}=0.
        \label{eq:asymex_fib_one_off}
    \end{align}
    The stacked square has bond dimension nine:
    \begin{align}
        B^{x'x} & =\sum_{m=0}^1 B_\tau^{x'm}\otimes B_\tau^{mx}.
        \label{eq:asymex_fib_stack}
    \end{align}
    Viewed as MPS tensors over the four-letter pair alphabet, its two
    target blocks are $B_1$ and $B_\tau$, each with multiplicity one
    and weight one.
\end{definition}

The review writes the operator with the prefactor $1/\sqrt{d_Ad_D}$
arising from the normalization of the fusion basis, where $d_A$ and
$d_D$ are the quantum dimensions of the labels on the adjacent
bonds~\cite[Appendix~A, ``The MPO for the Fibonacci model'']{Cirac2021Matrix};
that prefactor is omitted in the blocks defined here.

\begin{theorem}[The one-site matrices of the $\tau$ block]
    \label{thm:asymex_fib_entries}
    \lean{FibonacciCompression.fibTauMPS_apply_one,
        FibonacciCompression.fibTauMPS_apply_three}
    \leanok
    \uses{def:asymex_fib}
    The matrix of $B_\tau$ at the letter $(0,1)$ has the first row of
    (\ref{eq:asymex_fib_vertex}) as its first row and zero elsewhere, and
    its matrix at the letter $(1,1)$ has the third row of
    (\ref{eq:asymex_fib_vertex}) as its third row and zero elsewhere, with
    $\varphi^{-1/2}=\sqrt{(\sqrt5-1)/2}$.
\end{theorem}

\begin{proof}
    \leanok
    \uses{thm:asymex_golden_ring}
    Both matrices are obtained by embedding their coefficients in
    $\Z[\sigma]$ into $\C$ entrywise. Substitution of
    $\sigma=\sqrt{(\sqrt5-1)/2}$ gives
    (\ref{eq:asymex_fib_tau_01}) and
    (\ref{eq:asymex_fib_tau_11}).
\end{proof}

\begin{theorem}[Normality of the two Fibonacci blocks]
    \label{thm:asymex_fib_normal}
    \lean{FibonacciCompression.fibOneMPS_isNormal,
        FibonacciCompression.fibTauMPS_isNormal}
    \leanok
    \uses{def:asymex_fib, def:normal}
    Both blocks $B_1$ and $B_\tau$ are normal: their length-three words
    satisfy $\spn_\C\{B_1^w : |w|=3\} = \MN{2}$ and
    $\spn_\C\{B_\tau^w : |w|=3\} = \MN{3}$.
\end{theorem}

\begin{proof}
    \leanok
    \uses{thm:asymex_golden_ring}
    For the $\tau$ block, let $e_i$ be the $i$-th coordinate column
    vector and $f_i$ the $i$-th row of
    (\ref{eq:asymex_fib_vertex}). In this proof, superscripts
    $0,1,2,3$ enumerate the ordered pair alphabet. Thus
    $B_\tau^{i+1}=e_i f_i$ for $i\in\{0,1,2\}$, and
    \begin{align}
        B_\tau^{i+1}B_\tau^3B_\tau^{k+1}
         & =F_{i2}F_{2k}\,e_i f_k.
    \end{align}
    All factors $F_{i2}F_{2k}$ are nonzero, and the rows $f_k$
    form a basis. Each of the nine matrix units is an explicit
    combination over $\Z[\sigma]$ of these nine words.
    For the vacuum block, four length-three words suffice.
    Both families of coefficients are verified in $\Z[\sigma]$
    and transported to complex matrices by
    Theorem~\ref{thm:asymex_golden_ring}.
\end{proof}

\subsection{Admissible configurations and the vacuum operator}

Define the adjacency matrix $K=\begin{pmatrix}0&1\\1&1\end{pmatrix}$
and, for a cyclic word $x\in\{0,1\}^N$, set
\begin{align}
    p_N(x) & =\prod_{v=0}^{N-1}K_{x_vx_{v+1}}.
    \label{eq:asymex_fib_admissibility}
\end{align}
A word is admissible if $p_N(x)=1$, equivalently if it has no cyclic
neighboring pair $00$. This includes the last-to-first pair.
At $N=1$, the condition excludes the single label $0$.

\begin{lemma}[Vacuum projection and its kernel]
    \label{lem:asymex_fib_vacuum}
    For $N\geq1$, the vacuum operator $P_N=O_N(B_1)$ has matrix elements
    \begin{align}
        \bra{y}P_N\ket{x} & =\delta_{yx}p_N(x).
        \label{eq:asymex_fib_vacuum_kernel}
    \end{align}
    Thus $P_N\ket{x}=p_N(x)\ket{x}$. It is the orthogonal projection
    onto the admissible subspace:
    \begin{align}
        \ran P_N & =\spn_{\C}\{\ket{x}:p_N(x)=1\}.
        \label{eq:asymex_fib_vacuum_range}
    \end{align}
    \begin{align}
        \ker P_N & =\spn_{\C}\{\ket{x}:p_N(x)=0\}.
        \label{eq:asymex_fib_vacuum_null}
    \end{align}
    In particular, $P_N^2=P_N=P_N^\dagger$, but $P_N\neq\Id_{\mc H_N}$.
\end{lemma}

\begin{proof}
    The vacuum letters satisfy
    $(B_1^{yx})_{ab}=\delta_{yx}\delta_{ax}K_{xb}$.
    Expanding the periodic trace over bond labels
    $g=(g_0,\ldots,g_{N-1})$, with $g_N=g_0$, gives
    \begin{align}
        \bra{y}P_N\ket{x}
         & =\sum_g\prod_{v=0}^{N-1}
        \delta_{y_vx_v}\delta_{g_vx_v}K_{x_vg_{v+1}}
        =\delta_{yx}\prod_{v=0}^{N-1}K_{x_vx_{v+1}}.
    \end{align}
    Only $g=x$ contributes. Since $p_N(x)\in\{0,1\}$,
    the operator is diagonal with real idempotent entries.
    This proves the projection, range, and kernel statements.
    The nonzero vector $\ket{0,\ldots,0}$ lies in its kernel,
    so it is not the ambient identity. At one site,
    $P_1=\diag(0,1)$.
\end{proof}

\begin{theorem}[Vacuum unit laws]
    \label{thm:asymex_fib_unit}
    \lean{FibonacciCompression.fibOne_mul_fibOne,
        FibonacciCompression.fibOne_mul_fibTau,
        FibonacciCompression.fibTau_mul_fibOne}
    \leanok
    \uses{def:asymex_fib, def:mpo_family}
    For every $N\geq1$, the periodic operators satisfy
    \begin{align}
        O_N(B_1)^2 & =O_N(B_1).
        \label{eq:asymex_fib_unit_square}
    \end{align}
    \begin{align}
        O_N(B_1)O_N(B_\tau) & =O_N(B_\tau).
        \label{eq:asymex_fib_unit_left}
    \end{align}
    \begin{align}
        O_N(B_\tau)O_N(B_1) & =O_N(B_\tau).
        \label{eq:asymex_fib_unit_right}
    \end{align}
\end{theorem}

\begin{proof}
    \leanok
    For MPO tensors $T$ and $U$, denote their stacked product by
    $(TU)^{yx}=\sum_mT^{ym}\otimes U^{mx}$.
    Write $0_r$ for the tensor whose letters are zero $r\times r$
    matrices, and $\sim$ for simultaneous similarity of all letters.
    The three stacked products admit the decompositions
    \begin{align}
        B_1B_1 & \sim B_1\oplus0_2.
        \label{eq:asymex_fib_unit_blocks}
    \end{align}
    \begin{align}
        B_1B_\tau & \sim B_\tau\oplus0_3.
    \end{align}
    \begin{align}
        B_\tau B_1 & \sim B_\tau\oplus0_3.
    \end{align}
    The changes of bond coordinates and their inverses have entries
    in $\Z[\sigma]$. The decompositions, of dimensions $4=2+2$ and
    $6=3+3$, are verified by matrix multiplication with coefficients
    reduced using $\sigma^4=1-\sigma^2$.
    The periodic product formula for a multi-block compression gives
    the trace of the single nonzero target block. Since $N>0$, every
    zero block contributes zero, proving the three identities.
\end{proof}

By~(\ref{eq:asymex_fib_unit_left}), the range of $O_N(B_\tau)$ lies in
the admissible subspace. By~(\ref{eq:asymex_fib_unit_right}), it vanishes
on the orthogonal complement of that subspace. Thus the vacuum acts
as the identity on the admissible subspace and is the unit of the
algebra generated there by $O_N(B_\tau)$.

This vacuum block is distinct from the weighted string-net projector
of~\cite[Appendix~D.1]{Bultinck2017Anyons}, which combines the vacuum and
$\tau$ blocks with weights $w_1=1/(1+\varphi^2)$ and
$w_\tau=\varphi/(1+\varphi^2)$.

\subsection{Two periodic states and the Fibonacci action}

For an MPS tensor $D$, write
$\ket{V^{(N)}(D)}=\sum_x\tr(D^{x_0}\cdots D^{x_{N-1}})\ket{x}$.
These vectors are not normalized unless this is stated explicitly.

\begin{definition}[The all-$\tau$ and chain tensors]
    \label{def:asymex_fib_states}
    \lean{FibonacciCompression.fibAllTau,
        FibonacciCompression.fibChain}
    \leanok
    The all-$\tau$ tensor $A$ has bond dimension one, with
    $A^0=(0)$ and
    \begin{align}
        A^1 & =(1).
        \label{eq:asymex_fib_product_tensor}
    \end{align}
    The chain tensor $C$ has bond dimension two:
    \begin{align}
        C^0 & =\begin{pmatrix}0&\sigma\\0&0\end{pmatrix}.
        \label{eq:asymex_fib_chain_zero}
    \end{align}
    \begin{align}
        C^1 & =\begin{pmatrix}0&0\\1&-\sigma^2\end{pmatrix}.
        \label{eq:asymex_fib_chain_one}
    \end{align}
\end{definition}

\begin{lemma}[Amplitudes and norms]
    \label{lem:asymex_fib_amplitudes}
    Let $N\geq1$. The product tensor gives
    $\ket{V^{(N)}(A)}=\ket{1,\ldots,1}$, of norm one.
    For a word $x$, let $n_0(x)$ count its zeros and let $n_{11}(x)$
    count its cyclic neighboring pairs $11$. Then
    \begin{align}
        V^{(N)}(C)_x
         & =p_N(x)\sigma^{n_0(x)}(-\sigma^2)^{n_{11}(x)}.
        \label{eq:asymex_fib_chain_amplitude}
    \end{align}
    Both vectors are fixed by $P_N$. The chain norm is
    \begin{align}
        \bigl\|\ket{V^{(N)}(C)}\bigr\|^2
         & =\sum_{x\in\{0,1\}^N}
        p_N(x)\sigma^{2n_0(x)+4n_{11}(x)}.
        \label{eq:asymex_fib_chain_norm}
    \end{align}
    This norm is positive.
\end{lemma}

\begin{proof}
    The bond-one trace is $\prod_v\delta_{x_v1}$.
    For the chain, put
    $H=\begin{pmatrix}0&\sigma\\1&-\sigma^2\end{pmatrix}$.
    Since $(C^x)_{ab}=\delta_{ax}H_{xb}$, the cyclic contraction gives
    \begin{align}
        V^{(N)}(C)_x
         & =\sum_g\prod_v\delta_{g_vx_v}H_{x_vg_{v+1}}
        =\prod_vH_{x_vx_{v+1}}.
    \end{align}
    A pair $00$ makes the product zero. In an admissible word, every
    zero is followed by a one, giving $n_0(x)$ factors of $\sigma$.
    Each pair $10$ contributes $1$, and each pair $11$ contributes
    $-\sigma^2$. This proves~(\ref{eq:asymex_fib_chain_amplitude}).
    Both vectors have admissible support, so
    Lemma~\ref{lem:asymex_fib_vacuum} shows that $P_N$ fixes them.
    Summing squared amplitudes proves~(\ref{eq:asymex_fib_chain_norm}).
    The all-one amplitude is $(-\sigma^2)^N\neq0$.
\end{proof}

The normalized chain vector is obtained by dividing by the positive
square root of~(\ref{eq:asymex_fib_chain_norm}). For example,
$\ket{V^{(1)}(C)}=-\sigma^2\ket{1}$ and its squared norm is $\sigma^4$.
In particular, the two periodic vectors are not independent at every
length.

\begin{theorem}[Normality of the two state tensors]
    \label{thm:asymex_fib_states_normal}
    \lean{FibonacciCompression.fibAllTau_isNormal,
        FibonacciCompression.fibChain_isNormal}
    \leanok
    \uses{def:asymex_fib_states, def:normal}
    The length-one words of $A$ span $\MN{1}$, and the length-three
    words of $C$ span $\MN{2}$. Thus both tensors are normal.
\end{theorem}

\begin{proof}
    \leanok
    The letter $A^1$ spans $\MN{1}$. For $C$, write $E_{ab}$ for
    the matrix units of $\MN{2}$. Direct multiplication gives
    \begin{align}
        C^0C^1C^0 & =\sigma^2E_{01}.
    \end{align}
    \begin{align}
        C^0C^1C^1 & =-\sigma^3E_{00}+\sigma^5E_{01}.
    \end{align}
    \begin{align}
        C^1C^0C^1 & =\sigma C^1.
    \end{align}
    \begin{align}
        C^1C^1C^0 & =-\sigma^3E_{11}.
    \end{align}
    These four words span all four matrix units, since $\sigma\neq0$.
    The coefficients expressing the units in their span can all be
    chosen in $\Z[\sigma]$, using
    $\sigma^{-1}=\sigma+\sigma^3$.
    Embedding these identities into $\C$ proves that the words span the
    matrix algebra.
\end{proof}

\begin{theorem}[Action on the two periodic states]
    \label{thm:asymex_fib_state_action}
    \lean{FibonacciCompression.mpo_fibOne_mulVec_allTau,
        FibonacciCompression.mpo_fibTau_mulVec_allTau,
        FibonacciCompression.mpo_fibTau_mulVec_chain}
    \leanok
    \uses{def:asymex_fib, def:asymex_fib_states, def:mpo_family}
    For every $N\geq1$,
    \begin{align}
        O_N(B_1)\ket{V^{(N)}(A)}
         & =\ket{V^{(N)}(A)}.
        \label{eq:asymex_fib_action_unit}
    \end{align}
    \begin{align}
        O_N(B_\tau)\ket{V^{(N)}(A)}
         & =\ket{V^{(N)}(C)}.
        \label{eq:asymex_fib_action_product}
    \end{align}
    \begin{align}
        O_N(B_\tau)\ket{V^{(N)}(C)}
         & =\ket{V^{(N)}(A)}+\ket{V^{(N)}(C)}.
        \label{eq:asymex_fib_action_chain}
    \end{align}
\end{theorem}

\begin{proof}
    \leanok
    For an MPO tensor $T$ and MPS tensor $D$, their action tensor
    has letters $(T\cdot D)^y=\sum_xT^{yx}\otimes D^x$.
    As in the proof of Theorem~\ref{thm:asymex_fib_unit}, let $\sim$
    denote simultaneous similarity of letters and let $0_r$ denote
    a zero tensor of bond dimension $r$.
    The three action tensors have block decompositions
    \begin{align}
        B_1\cdot A & \sim A\oplus0_1.
        \label{eq:asymex_fib_action_blocks_unit}
    \end{align}
    \begin{align}
        B_\tau\cdot A & \sim C\oplus0_1.
        \label{eq:asymex_fib_action_blocks_product}
    \end{align}
    \begin{align}
        B_\tau\cdot C & \sim A\oplus C\oplus0_3.
        \label{eq:asymex_fib_action_blocks_chain}
    \end{align}
    In lexicographic order on the product bond indices, matrices whose
    columns give these adapted bases are respectively
    \begin{align}
        Q_2 & =\begin{pmatrix}0&1\\1&-1\end{pmatrix}.
    \end{align}
    \begin{align}
        Q_3 & =\begin{pmatrix}1&0&1\\0&0&-\sigma\\0&1&\sigma^2\end{pmatrix}.
    \end{align}
    \begin{align}
        Q_6 & =\begin{pmatrix}
                   0               & 0 & 0         & 0          & 1        & 1         \\
                   0               & 1 & 0         & 0          & 0        & \sigma^2  \\
                   1               & 0 & \sigma    & 1          & -\sigma  & -\sigma^3 \\
                   0               & 0 & 0         & 1+\sigma^2 & 0        & 0         \\
                   0               & 0 & 0         & 0          & \sigma^2 & 0         \\
                   \sigma+\sigma^3 & 0 & -\sigma^2 & 0          & 0        & -\sigma^2
               \end{pmatrix}.
    \end{align}
    They are invertible, with inverses over $\Z[\sigma]$.
    For each action tensor, let $D_{\mathrm{blk}}$ denote the
    corresponding block tensor in
    (\ref{eq:asymex_fib_action_blocks_unit}),
    (\ref{eq:asymex_fib_action_blocks_product}), or
    (\ref{eq:asymex_fib_action_blocks_chain}).
    Matrix multiplication verifies
    $(T\cdot D)^yQ=QD_{\mathrm{blk}}^y$.
    All off-diagonal blocks vanish, giving the three multi-block
    compressions. The periodic action formula for each compression
    sums its target vectors. The zero blocks contribute nothing
    because $N>0$. The first two sums have one target each, and the
    third has the targets $A$ and $C$.
\end{proof}

The first nontrivial action also has a direct contraction proof.
The all-one input selects the letters $B_\tau^{01}$ and
$B_\tau^{11}$. Both have zero row $1$, so a cyclic bond configuration
using that coordinate contributes zero. Restricting to the ordered
coordinates $(0,2)$ gives exactly $C^0$ and $C^1$. Hence
\begin{align}
    \bra{x}O_N(B_\tau)\ket{1,\ldots,1}
     & =\tr(B_\tau^{x_0 1}\cdots B_\tau^{x_{N-1}1})
    =\tr(C^{x_0}\cdots C^{x_{N-1}}).
\end{align}
There is no additional normalization factor.

Relative to the ordered pair of these unnormalized state families,
the action coefficients are the Fibonacci matrix
$N_\tau=\begin{pmatrix}0&1\\1&1\end{pmatrix}$, realizing the nonnegative-integer
action discussed in~\cite{GarreRubio2023MPOSym}.
The positive-length hypothesis is essential: at $N=0$ the periodic trace is the
bond dimension, yielding $O_0(B_1)=2$, $O_0(B_\tau)=3$, and empty amplitudes
$1$ and $2$ for $A$ and $C$.

\subsection{Asymmetric compression of the stacked square}

The stacked square in~(\ref{eq:asymex_fib_stack}) contracts the physical
index between two copies of $B_\tau$ while retaining both virtual
indices. Its periodic operator is therefore the square of
$O_N(B_\tau)$. The compression maps onto its two target blocks also
have entries in $\Z[\sigma]$, so the fusion rule below is again an
exact computation in the ring of golden integers.

\begin{center}
    % Formula: B^{x'x}=\sum_m(B_\tau)^{x'm}\otimes(B_\tau)^{mx}, the
    % stacked product of Definition~\ref{def:asymex_fib}, realizing
    % $\tau\otimes\tau$; Theorem~\ref{thm:asymex_fib_fusion} identifies
    % its periodic trace with that of $B_1\oplus B_\tau$, the
    % realization $1\oplus\tau$.
    % Source: Definition~\ref{def:asymex_fib}.
    % Ink-to-index map: both rows are the $\tau$ tensor $B_\tau$; the
    % shared vertical pairleg at each site is the summed intermediate
    % label $m$; the drawn upper-row label $y$ is the per-site
    % component of the outgoing index $x'$.
    % Contraction set: the vertical pairlegs (index $m$) at every site
    % and the horizontal virtual bonds within each row; the physical
    % legs $(y,x)$, the per-site components of $(x',x)$, remain open.
    % Boundary signature: (virtual-west=0 | virtual-east=0 |
    % physical-up=3 | physical-down=3) on both sides.
    \tenkzkernel
    \begin{tenkz}[rows={op,op}, cols=4, boundary=periodic]
        \tn[skin=mpo, ports={90:physical:$y_1$, 270:physical}]{B_\tau} &
        \tn[skin=mpo, ports={90:physical:$y_2$, 270:physical}]{B_\tau} &
        \tn[skin=dots, ports={180:virtual, 0:virtual}]{} &
        \tn[skin=mpo, ports={90:physical:$y_N$, 270:physical}]{B_\tau} \\
        \tn[skin=mpo, ports={90:physical, 270:physical:$x_1$}]{B_\tau} &
        \tn[skin=mpo, ports={90:physical, 270:physical:$x_2$}]{B_\tau} &
        \tn[skin=dots, ports={180:virtual, 0:virtual}]{} &
        \tn[skin=mpo, ports={90:physical, 270:physical:$x_N$}]{B_\tau}
    \end{tenkz}
    \quad $=$ \quad
    \begin{tenkz}[rows={op}, cols=4, boundary=periodic]
        \tn[skin=mpo, ports={90:physical:$y_1$, 270:physical:$x_1$}]{B} &
        \tn[skin=mpo, ports={90:physical:$y_2$, 270:physical:$x_2$}]{B} &
        \tn[skin=dots, ports={180:virtual, 0:virtual}]{} &
        \tn[skin=mpo, ports={90:physical:$y_N$, 270:physical:$x_N$}]{B}
    \end{tenkz}
\end{center}

\begin{theorem}[Split compression of the Fibonacci square]
    \label{thm:asymex_fib_compression}
    % Principal declaration: FibonacciCompression.fibonacci_isReduction.
    \lean{FibonacciCompression.fibonacci_compression,
        FibonacciCompression.fibonacci_isReduction,
        FibonacciCompression.fibonacci_left_mul_right_of_ne,
        FibonacciCompression.fibonacci_dim_eq,
        FibonacciCompression.fibonacci_evalWord_remainder_eq_zero,
        FibonacciCompression.fibonacci_remainder_eq_zero}
    \leanok
    \uses{def:asymex_fib, thm:asym_multiblock}
    The source of Definition~\ref{def:asymex_fib} carries the block
    triangular form of Theorem~\ref{thm:asym_multiblock}, with target
    blocks $C_1=B_1$ and $C_\tau=B_\tau$, and $z=4$ zero slots.
    The dimension count is $9=2+3+4$. Here the indexed tensors $C_s$
    denote the MPO target blocks, not the chain tensor $C$.
    There are compression maps
    $V_s:\C^{D_s}\to\C^9$ and $W_s:\C^9\to\C^{D_s}$,
    where $D_1=2$ and $D_\tau=3$, such that
    $W_sV_s=\Id_{D_s}$ and $W_sV_t=0$ for $s\neq t$.
    For every word $w$ over the pair alphabet,
    \begin{align}
        W_sB^wV_s & =C_s^w.
        \label{eq:asymex_fib_reduction}
    \end{align}
    Every remainder word of length at least six vanishes.
    In fact, each remainder letter is zero, so the extension splits.
\end{theorem}

\begin{proof}
    \leanok
    \uses{thm:asymex_explicit_gauge, thm:asymex_golden_ring}
    The adapted basis has as its first columns the two fusion tensors,
    followed by a basis of the four-dimensional subspace annihilated
    by every source letter. Writing $Q$ for this basis matrix,
    the finite matrix calculation over $\Z[\sigma]$ gives
    \begin{align}
        Q^{-1}B^{x'x}Q
         & =B_1^{x'x}\oplus B_\tau^{x'x}\oplus0_4.
        \label{eq:asymex_fib_split_letters}
    \end{align}
    The corresponding column blocks of $Q$ and row blocks of $Q^{-1}$
    give $V_s$ and $W_s$. Their products give biorthogonality, and
    multiplication of~(\ref{eq:asymex_fib_split_letters}) along a word
    gives~(\ref{eq:asymex_fib_reduction}).
    Block diagonality supplies triangularity, the target diagonal
    blocks, and the vanishing remainder. The generic nilpotency
    bound of clause (vi) of Theorem~\ref{thm:asym_multiblock} is
    $|S|+z=6$, where $S$ is the two-element set of target blocks;
    the dimension count is clause (vii).
    The four zero directions are virtual bond directions, not
    physical configurations of the chain.
\end{proof}

\begin{theorem}[The fusion rule]
    \label{thm:asymex_fib_fusion}
    % Principal declaration: FibonacciCompression.fibonacci_fusion_rule.
    \lean{FibonacciCompression.fibStack_trace_evalWord,
        FibonacciCompression.fibonacci_fusion_rule,
        FibonacciCompression.fibStack_mpv}
    \leanok
    \uses{def:asymex_fib, def:mpo_family}
    For every nonempty word $w$ over the alphabet of label pairs,
    \begin{align}
        \tr(B^w) & =\tr(B_1^w)+\tr(B_\tau^w).
        \label{eq:asymex_fib_trace}
    \end{align}
    Thus at every positive system size $N$ the periodic operators
    satisfy~(\ref{eq:asymex_fib_fusion}), and the periodic vector
    of the stacked tensor is the sum of the periodic vectors of
    the two target blocks.
\end{theorem}

\begin{proof}
    \leanok
    \uses{thm:asymex_fib_compression, prop:asym_compression_trace}
    Apply Proposition~\ref{prop:asym_compression_trace} to
    Theorem~\ref{thm:asymex_fib_compression}. This gives
    (\ref{eq:asymex_fib_trace}).
    For output and input words $y,x$, the periodic product contracts
    as
    \begin{align}
        \bra{y}O_N(B_\tau)^2\ket{x}
         & =\sum_{m\in\{0,1\}^N}
        \tr(B_\tau^{y_0m_0}\cdots B_\tau^{y_{N-1}m_{N-1}})
        \tr(B_\tau^{m_0x_0}\cdots B_\tau^{m_{N-1}x_{N-1}})
        \nonumber                                     \\
         & =\tr(B^{y_0x_0}\cdots B^{y_{N-1}x_{N-1}}).
    \end{align}
    Substitution of~(\ref{eq:asymex_fib_trace}) proves the operator
    identity entry by entry. The periodic-vector statement is the
    same trace identity read over the pair alphabet.
\end{proof}

At length zero, the traces of the stacked tensor and the two targets
are $9$, $2$, and $3$, respectively. The omitted four-dimensional
zero block has empty trace four, although all its positive-length
traces vanish. This explains the positive-length restriction in the
fusion rule without changing the trace convention.

\begin{theorem}[The fusion tensors]
    \label{thm:asymex_fib_fusion_tensors}
    % Principal declaration: FibonacciCompression.fibStack_mul_fusionRight.
    \lean{FibonacciCompression.fibFusionRight, FibonacciCompression.fibFusionLeft,
        FibonacciCompression.fibonacci_right_eq, FibonacciCompression.fibonacci_left_eq,
        FibonacciCompression.fibStack_mul_fusionRight,
        FibonacciCompression.fusionLeft_mul_fibStack}
    \leanok
    \uses{def:asymex_fib, thm:asymex_fib_compression}
    The compression maps of the two target blocks are the fusion tensors
    of the Fibonacci algebra. The map $V_s$ is the column block of the
    adapted basis carrying $s$, and $W_s$ is the matching row block of
    its inverse. For every physical letter and $s\in\{1,\tau\}$,
    \begin{align}
        B^{x'x}V_s & =V_sC_s^{x'x}.
        \label{eq:asymex_fib_right_intertwiner}
    \end{align}
    \begin{align}
        W_sB^{x'x} & =C_s^{x'x}W_s.
        \label{eq:asymex_fib_left_intertwiner}
    \end{align}
\end{theorem}

\begin{proof}
    \leanok
    \uses{thm:asymex_explicit_gauge}
    Theorem~\ref{thm:asymex_explicit_gauge} identifies the maps with
    the row and column blocks of the adapted basis.
    Each intertwining relation is verified by multiplication over
    $\Z[\sigma]$ and transported to complex matrices by
    Theorem~\ref{thm:asymex_golden_ring}.
    Equivalently, taking the corresponding column or row block
    of~(\ref{eq:asymex_fib_split_letters}) gives
    (\ref{eq:asymex_fib_right_intertwiner}) or
    (\ref{eq:asymex_fib_left_intertwiner}).
\end{proof}

\begin{center}
    % Formula: B^{x'x}V_s=V_sC_s^{x'x} and W_sB^{x'x}=C_s^{x'x}W_s of
    % Theorem~\ref{thm:asymex_fib_fusion_tensors}, for $s\in\{1,\tau\}$.
    % Source: Theorem~\ref{thm:asymex_fib_fusion_tensors}.
    % Ink-to-index map: the ring is the fusion tensor $V_s$ or $W_s$;
    % the square node is the stacked source $B$ or the target block
    % $C_s\in\{B_1,B_\tau\}$.
    % Contraction set: the single virtual bond joining the ring to the
    % square node in each pair; the physical pair leg $(x',x)$ remains
    % open.
    % Boundary signature: (virtual-west=1 | virtual-east=1 |
    % physical-up=1 | physical-down=0) for each pictured equality.
    \tenkzkernel
    \begin{tenkz}[rows={wire}, cols=2, west=open, east=open]
        \tn[label pos=s, ports={90:physical:$x'x$}]{B} & \tn[skin=ring]{V_s}
    \end{tenkz}
    \quad $=$ \quad
    \begin{tenkz}[rows={wire}, cols=2, west=open, east=open]
        \tn[skin=ring]{V_s} & \tn[label pos=s, ports={90:physical:$x'x$}]{C_s}
    \end{tenkz}
    \\
    \begin{tenkz}[rows={wire}, cols=2, west=open, east=open]
        \tn[skin=ring]{W_s} & \tn[label pos=s, ports={90:physical:$x'x$}]{B}
    \end{tenkz}
    \quad $=$ \quad
    \begin{tenkz}[rows={wire}, cols=2, west=open, east=open]
        \tn[label pos=s, ports={90:physical:$x'x$}]{C_s} & \tn[skin=ring]{W_s}
    \end{tenkz}
\end{center}
